% Generated draft: requires mathematical and source review.
\documentclass[11pt]{article}
\usepackage[margin=25mm]{geometry}
\usepackage{fontspec}
\setmainfont{Latin Modern Roman}
\usepackage{amsmath,amssymb,mathtools,unicode-math}
\setmathfont{Latin Modern Math}
\usepackage{enumitem,needspace,xurl,hyperref}
\setlength{\parindent}{0pt}
\setlength{\parskip}{5pt}
\setlength{\emergencystretch}{3em}
\newcommand{\N}{\mathbb{N}}
\newcommand{\Z}{\mathbb{Z}}
\newcommand{\Q}{\mathbb{Q}}
\newcommand{\R}{\mathbb{R}}
\newcommand{\C}{\mathbb{C}}
\newcommand{\F}{\mathbb{F}}
\newcommand{\A}{\mathbb{A}}
\newcommand{\PP}{\mathbb{P}}
\newcommand{\E}{\mathbb{E}}
\newcommand{\eps}{\varepsilon}
\newcommand{\dd}{\,\mathrm{d}}
\begin{document}
\section*{UnsolvedMath ID: 30005652}
\textbf{Source problem code:} catalogue entry
\subsection{Definitions and mathematical statement}
Let \(k\) be a field of characteristic \(p>0\), and let \(G\) be a finite flat group scheme over \(k\). In the finite-dimensional convention, the stable module category of \(G\) has finite-dimensional \(G\)-representations as objects and morphisms modulo those factoring through projective representations. The principal block is the block containing the trivial representation \(k\). The supplied record does not specify whether finite-dimensional representations or a larger stable module category are intended.

Write \(\mathcal C_0\) for the principal block of an \(\infty\)-categorical enhancement intended by the source. To define the requested invariants, one needs a specified symmetric monoidal structure and suitable categorical hypotheses. When these data are available, \(\operatorname{Pic}(\mathcal C_0)\) is the group of equivalence classes of tensor-invertible objects, with multiplication induced by tensor product. The categorical Brauer group \(\operatorname{Br}(\mathcal C_0)\) is the group of Morita-equivalence classes of Azumaya algebra objects, under tensor product, using the chosen categorical convention. The record does not specify that convention or explain how the principal block acquires the required monoidal structure; a block decomposition alone does not supply this information.

\textbf{Source question, conditional on the missing categorical data.} For every such \(k\) and \(G\), and for the enhancements and principal blocks covered by the source's intended hypotheses, can descent spectral sequences based on the associated spectral scheme \(X\) completely determine both
\[
\operatorname{Pic}(\mathcal C_0)
\qquad\text{and}\qquad
\operatorname{Br}(\mathcal C_0)?
\]
Here a complete computation must determine the groups themselves, including any differentials, convergence issues, and extension problems needed to recover them from the spectral sequences. The supplied text does not identify \(X\), the asserted equivalence involving \(X\), the descent topology, the coefficient objects, or the precise spectral sequences. Consequently, it does not determine a fully specified universal mathematical assertion.

The original excerpt describes principal blocks of \emph{any} \(\infty\)-categorical enhancement of stable module categories of finite flat group schemes as examples satisfying omitted preceding hypotheses, attributing this to [FP07]. Those hypotheses and the referenced equivalences are absent from the record, so the scope of this assertion remains unresolved.

The excerpt also contains a distinct question, explicitly labelled \emph{Question 12}: do the spectral schemes in those omitted equivalences admit natural spectral moduli-theoretic interpretations? Its unidentified speaker reports an affirmative answer only for elementary abelian groups. This adjacent question and qualification are preserved as source context; the record does not establish whether Question 12 is the intended target of the registered entry.
\subsection{Short English statement}
The recorded main question asks whether descent from associated spectral schemes can completely compute the Picard and Brauer groups of principal stable-module blocks, including the intended examples from finite flat group schemes in positive characteristic. Essential categorical and descent data are missing. The original excerpt also includes a separate moduli-interpretation question, leaving the precise intended target uncertain.
\subsection{Sources}
\begin{itemize}
\item Supplied local UnsolvedMath registry record: permanent ID \(30005652\), title \emph{Picard and Brauer Groups of Stable Module Blocks}, exact source problem code \texttt{catalogue entry}. The repository registry locations are \texttt{data/unsolvedmath/problems.json} and \texttt{catalog/problems.json}; no independent inspection of either is asserted here. No proposer is named; the supplied proposed year is 2023.
\item Supplied report citation: \emph{Tensor-Triangular Geometry and Interactions} (2024), Oberwolfach Reports, pp.\ 2303--2358, DOI \texttt{10.4171/owr/2023/40}. \url{https://doi.org/10.4171/owr/2023/40}. The background describes Oberwolfach Reports as documenting mathematical workshops and their open-problem discussions. No particular contribution, speaker, or problem page is identified.
\item Supplied literature source, labelled ``Hill--Hopkins, stable Picard groups via descent'': \url{https://arxiv.org/abs/1601.03098}. The attribution and claimed relevance are supplied metadata and have not been independently verified.
\item Supplied literature source, labelled ``OWR report posing the principal-block question'': \url{https://doi.org/10.4171/owr/2023/40}.
\item Supplied literature source, labelled ``Krause, Picard groups in equivariant homotopy theory'': \url{https://doi.org/10.1112/topo.70020}. The attribution and claimed relevance are supplied metadata and have not been independently verified.
\item The original excerpt cites [FP07] for the finite-flat-group-scheme examples, [MS16] for an \(\mathrm{TMF}\) Picard-group computation using an \'etale descent spectral sequence, and [BMS22] for a classification of certain Azumaya algebras for \(\mathrm{TMF}\). The record supplies no bibliographic entries identifying these labels; no identifications are inferred.
\item Literature assessment dated exactly \texttt{2026-08-22}: the supplied curation labels the problem \emph{partially solved} and states that descent spectral sequences compute Picard and Brauer invariants in important stable-module and chromatic examples, while a complete computation for principal blocks of arbitrary finite flat group schemes is not established. This is supplied dated curation, explicitly described in the background as literature triage requiring independent verification; it is not an independently verified assessment here.
\item The registry labels the statement \texttt{corrected\_verified} and records a claim that the official report was checked directly and the clean statement preserves the complete target. These are supplied verification claims, not verification performed for this draft; the missing context and adjacent Question 12 remain unresolved.
\end{itemize}
\end{document}
